\documentclass[11pt, UTF8]{ctexart}
\usepackage{amsmath, amsfonts, amssymb}
\usepackage{extarrows} 
\usepackage{geometry} 
\usepackage{xcolor} 
\usepackage{enumitem} 
\usepackage[most]{tcolorbox} 
\usepackage{fancyhdr} 
\usepackage{diagbox} 
\usepackage{tikz}    
\setlength{\headheight}{13.6pt}
\geometry{a4paper, margin=0.8in} 
\pagestyle{fancy}
\fancyhf{}
\lhead{\color{fblue} \hwdate 作业 答案} 
\rhead{\color{fblue} \today}
\cfoot{\thepage}

\definecolor{fblue}{RGB}{0, 74, 142}
\definecolor{fred}{RGB}{204, 26, 26}

\newtcolorbox{problem}[2]{
    enhanced,
    colback=white,
    colframe=fblue!80,
    arc=2mm,
    boxrule=0.8pt,
    before upper={\textbf{\color{fblue} 问题 #1 (#2)}\hspace{0.5em}},
    before skip=1.5em,
    after skip=1.5em,
    left=3mm, right=3mm, top=2mm, bottom=2mm
}

\newcommand{\hwdate}{4月30日} % 可根据需要修改日期
\newcommand{\Prob}{P}
\newcommand{\comp}[1]{\overline{#1}\mskip 3mu}
\newcommand{\ans}[1]{\mathbf{\color{fred}#1}}

\begin{document}

\begin{center}
    {\LARGE \bfseries \color{fblue} \hwdate 作业 答案}
\end{center}

% ---------------- 题目 1 ----------------
\begin{problem}{1}{3.26}
设随机变量 $X,Y$ 相互独立，它们的概率密度均为
\[ f(x) = \begin{cases} \mathrm{e}^{-x}, & x > 0, \\ 0, & \text{其他.} \end{cases} \]
求 $Z = \frac{Y}{X}$ 的概率密度.
\end{problem}

\paragraph{\color{fblue}解} 
\[ f_X(x) = \begin{cases} \mathrm{e}^{-x}, & x > 0, \\ 0, & \text{其他,} \end{cases} \qquad f_Y(y) = \begin{cases} \mathrm{e}^{-y}, & y > 0, \\ 0, & \text{其他.} \end{cases} \]
由公式
\[ f_Z(z) = \int_{-\infty}^{\infty} |x| f_X(x)f_Y(xz)\mathrm{d}x, \]
仅当 $\begin{cases} x > 0, \\ xz > 0, \end{cases}$ 即 $\begin{cases} x > 0, \\ z > 0 \end{cases}$ 时，上述积分的被积函数不等于零，于是当 $z > 0$ 时有
\[ f_Z(z) = \int_{0}^{\infty} x\mathrm{e}^{-x}\mathrm{e}^{-xz}\mathrm{d}x = \int_{0}^{\infty} x\mathrm{e}^{-x(z+1)}\mathrm{d}x = \ans{\frac{1}{(z+1)^2}}. \]
当 $z \leqslant 0$ 时 $f_Z(z) = 0$, 即
\[ f_Z(z) = \begin{cases} \ans{\frac{1}{(z+1)^2},} & \ans{z > 0,} \\ \ans{0,} & \ans{z \leqslant 0.} \end{cases} \]

\vspace{1.5em}

% ---------------- 题目 2 ----------------
\begin{problem}{2}{3.28}
设 $X,Y$ 是相互独立的随机变量，它们都服从正态分布 $N(0,\sigma^2)$. 试验证随机变量 $Z = \sqrt{X^2+Y^2}$ 的概率密度为
\[ f_Z(z) = \begin{cases} \frac{z}{\sigma^2}\mathrm{e}^{-z^2/(2\sigma^2)}, & z > 0, \\ 0, & \text{其他.} \end{cases} \]
我们称 $Z$ 服从参数为 $\sigma$ $(\sigma > 0)$ 的瑞利 (Rayleigh) 分布.
\end{problem}

\paragraph{\color{fblue}证}
先来求 $Z$ 的分布函数 $F_Z(z)$, 由于 $Z = \sqrt{X^2+Y^2} \geqslant 0$, 知当 $z < 0$ 时， $F_Z(z) = 0$. 当 $z \geqslant 0$ 时有
\begin{align*}
F_Z(z) &= P\{Z \leqslant z\} = P\{\sqrt{X^2+Y^2} \leqslant z\} \\
&= P\{X^2+Y^2 \leqslant z^2\} = \iint\limits_{x^2+y^2 \leqslant z^2} f(x,y)\mathrm{d}x\mathrm{d}y \\
&= \iint\limits_{x^2+y^2 \leqslant z^2} \frac{1}{2\pi\sigma^2}\mathrm{e}^{-(x^2+y^2)/(2\sigma^2)}\mathrm{d}x\mathrm{d}y \\
&\xlongequal{\text{用极坐标}} \int_{0}^{2\pi} \mathrm{d}\theta \int_{0}^{z} \frac{1}{2\pi\sigma^2}\mathrm{e}^{-r^2/(2\sigma^2)} r\mathrm{d}r \\
&= 2\pi \cdot \frac{1}{2\pi} \left[ -\mathrm{e}^{-r^2/(2\sigma^2)} \right] \Bigg|_{0}^{z} = \ans{1 - \mathrm{e}^{-z^2/(2\sigma^2)}}.
\end{align*}
将 $F_Z(z)$ 关于 $z$ 求导数，得 $Z$ 的概率密度为
\[ f_Z(z) = \begin{cases} \ans{\frac{z}{\sigma^2}\mathrm{e}^{-z^2/(2\sigma^2)},} & \ans{z > 0,} \\ \ans{0,} & \text{\textbf{\color{fred}其他.}} \end{cases} \]

\vspace{1.5em}

% ---------------- 题目 3 ----------------
\begin{problem}{3}{3.31}
对某种电子装置的输出测量了 $5$ 次，得到结果为 $X_1, X_2, X_3, X_4, X_5$. 设它们是相互独立的随机变量且都服从参数 $\sigma = 2$ 的瑞利分布. \\
(1) 求 $Z = \max\{X_1, X_2, X_3, X_4, X_5\}$ 的分布函数. \\
(2) 求 $P\{Z > 4\}$.
\end{problem}

\paragraph{\color{fblue}解}
参数 $\sigma = 2$ 的瑞利分布，其概率密度为
\[ f_X(x) = \begin{cases} \frac{x}{4}\mathrm{e}^{-x^2/8}, & x > 0, \\ 0, & \text{其他.} \end{cases} \]
设其分布函数为 $F_X(x)$. 则当 $x < 0$ 时，$F_X(x) = 0$, 当 $x \geqslant 0$ 时有
\[ F_X(x) = \int_{0}^{x} \frac{x}{4}\mathrm{e}^{-x^2/8}\mathrm{d}x = -\mathrm{e}^{-x^2/8} \Bigg|_{0}^{x} = 1 - \mathrm{e}^{-x^2/8}. \]
即有
\[ F_X(x) = \begin{cases} 1 - \mathrm{e}^{-x^2/8}, & x \geqslant 0, \\ 0, & x < 0. \end{cases} \]
(1) 因 $X_1, X_2, X_3, X_4, X_5$ 相互独立，且都服从参数 $\sigma = 2$ 的瑞利分布，故 $Z = \max\{X_1, X_2, X_3, X_4, X_5\}$ 的分布函数为
\[ F_Z(z) = [F_X(z)]^5 = \begin{cases} \ans{(1 - \mathrm{e}^{-z^2/8})^5,} & \ans{z \geqslant 0,} \\ \ans{0,} & \ans{z < 0.} \end{cases} \]

(2) 
\begin{align*} 
P\{Z > 4\} &= 1 - P\{Z \leqslant 4\} = 1 - F_Z(4) \\
&= 1 - (1 - \mathrm{e}^{-2})^5 = \ans{0.516\,7}.
\end{align*}

\end{document}